% \newpage


\section{Related Work}
% \paragraph{Accelerating Scientific Discovery with LLMs.}
% Recent work demonstrates that LLMs have the potential to break the scientific boundary built by human researchers. For example, AlphaEvolve \citep{novikov2025alphaevolve}, as an evolutionary coding agent, developed a search algorithm that found a procedure to multiply two 4×4 complex-valued matrices using 48 scalar multiplications, improves Strassen (1969)’s algorithm.
% ASI-ARCH \citep{liu2025alphago} discovers 106 state-of-the-art linear attention architectures by combining LLMs with neural architecture search (NAS). AlphaResearch also finds the best-of-known \emph{``Packing Circles''} solutions, improving the previous algorithm discovered by human and AlphaEvolve.
\paragraph{LLMs for New Ideas.}
Several recent works explored methods to improve research idea generation, such as iterative novelty refinement \citep{wang2024scimon, Baek2024ResearchAgentIR}. These works focus on improving the research idea over vanilla prompting but critically miss an effective verification method. To promote more reliable AI-generated research ideas, many studies have proposed solutions from different perspectives, such as comparisons with any human expert \citep{si2024can}, using LLMs for executing experiments by generating code with human-curated research problems \citep{Huang2023MLAgentBenchEL, Tian2024SciCodeAR}, and executing LLM-generated research ideas with LLM-generated programs \citep{Li2024MLRCopilotAM, lu2024ai, aygun2025ai}. These works either use automatic program evaluation or unverifiable LLM evaluator method, which presents a challenge for their scalability to real-world advanced algorithm discovery. 
Our AlphaResearch presents a more feasible direction by combining program execution with RM training from real-world peer-reviewed research records.

\paragraph{LLMs for Code Generation.}
In autonomous research agents, code generation serves as a fundamental step. Previous models~\citep{guo2024deepseek, yu2023wavecoder, hui2024qwen2} and benchmarks ~\citep{chen2021evaluating,yu-etal-2025-humaneval} for code generation are in a longstanding pursuit of synthesizing code from natural language descriptions. SWE-Bench \citep{jimenez2024swebench}, PaperBench~\cite{starace2025paperbench}, MLE-Bench~\cite{chan2024mle} introduces the problems in real-world agentic coding. Many studies on SWE-Bench have greatly contributed to the emergence of coding agents like SWE-Agent \citep{yang2024swe} and OpenHands \citep{wang2025openhands}. These agent frameworks greatly facilitate the training of agentic LLMs like Kimi-K2 \citep{team2025kimi} and GLM-4.5 \citep{zeng2025glm}.
The surge of these models on SWE-Bench underscores a critical need to reassess the future directions of coding agent research.
% Our AlphaResearchComp benchmark shows that testing LLMs on open-ended research for algorithm discovery is a promising direction to adapt language models to real-world tasks.



